% day1_seg1_heuristic/update.tex
\begin{frame}
  \onslide<1->{\frametitle{Measurement Update (Correction)}}
 \begin{itemize}
  \onslide<2->{\item In running average example, the update implicitly
 assumed that measurements were corrupted by white noise}
   \onslide<3->{\item Math model for this would be:\\ $z(k)=x(k)+v(k)$, \\where $v(k)$ is
  Gaussian white measurement noise sequence}
   \onslide<4->{\item Measurement processing is also often called Observation processing}
\end{itemize}
\end{frame}

\begin{frame}
  \onslide<1->{\frametitle{Measurement Update (Correction)}}
 \begin{itemize}
   \onslide<2->{\item Because we assumed that measurement noise was white
 Gaussian and zero mean, our best guess was simply the average value
 of all the measurements}
   \onslide<3->{\item Remember the resistor example: our voltmeter is not
 perfect, and each measurement will differ slightly from the previous
 one, even though we are measuring the same resistor each time}
\end{itemize}
\end{frame}

\begin{frame}
  \onslide<1->{\frametitle{Measurement Update (Correction)}}
 \begin{itemize}
 \onslide<2->{\item General Kalman filtering complicates the measurement model
 by adding two new elements:}
 \begin{itemize}
 \onslide<3->{\item The measurement may have a more complicated relationship to
 the underlying system state:  
\[z(k)=h\left[k,x(k)\right]+v(k),\] 
where $h$ could be a non-linear function with possible explicit time dependence}
 \onslide<4->{\item Just as the state $x$ is allowed to be a vector variable, so
 the measurement $z$ could also be a vector variable}
 \end{itemize}
\end{itemize}
\end{frame}
